\documentclass[runningheads]{llncs}
\usepackage[T1]{fontenc}
\usepackage[utf8]{inputenc}
\usepackage{lmodern}
\usepackage{amsmath,amssymb,graphicx,booktabs,array,xcolor}
\usepackage{listings}
\usepackage{hyperref}
\hypersetup{colorlinks=true,linkcolor=black,citecolor=black,urlcolor=blue,pdftitle={Integrating Native Prime-Field Reasoning into Z3},pdfauthor={Romain Soulat and Jean Frédéric Etienne}}
\lstset{basicstyle=\ttfamily\small,columns=fullflexible,keepspaces=true,breaklines=true,frame=single,framesep=5pt,showstringspaces=false}
\newcommand{\FF}{\mathbb{F}}
\newcommand{\tool}{Z3+FF}
\newcommand{\qfff}{\texttt{QF\_FF}}
\newcommand{\flag}[1]{\texttt{#1}}
\input{numbers.tex}
\input{round8-numbers.tex}
\begin{document}
\title{Integrating Native Prime-Field Reasoning into Z3}
\titlerunning{Native Prime-Field Reasoning in Z3}
\author{Romain Soulat \and Jean Frédéric Etienne}
\authorrunning{R. Soulat and J. F. Etienne}
\institute{Input Output Group}
\maketitle
\begin{abstract}
Finite-field constraints arise naturally in the verification of zero-knowledge
circuits, yet their use in a general-purpose SMT workflow requires more than a
standalone polynomial solver. We present \tool{}, an experimental integration
of prime-field reasoning into Z3. It provides native field syntax and models,
C/C++/Python interfaces, circuit-aware preprocessing, bounded modular algebra,
and ground combination with Z3's other theories. Candidate field models guide
shared-equality decisions; an exact, range-constrained bit-vector bridge supplies
a fallback without assuming stable infiniteness. The implementation introduces
no CoCoA dependency. A complete historical evaluation covers all \CorpusN{}
byte-distinct inputs from three public artifacts: the previous Z3 default solves
\ZthreeSolved{}, compared with \CvcSolved{} and \SplitSolved{} for cvc5 1.3.4
GB and split. A new paired study enables bounded root completion and diversified
model search, improving coverage from 886 to 888 on 1,256 inputs with no lost
solves and 0.45\% higher summed time on common successes. Category plots expose
remaining gaps and confirmed short-case slowdowns. The new default has not been
measured across the full corpus. Proof reconstruction and large-modulus
primality certificates remain outside the implemented scope.
\keywords{SMT \and Finite fields \and Z3 \and Theory combination \and Circuit verification}
\end{abstract}

\section{Introduction}

Arithmetic circuits used in zero-knowledge systems encode computation over a
finite field. A verification query may ask whether two circuit encodings agree,
whether outputs are determined by inputs, or whether arithmetic constraints
preserve a source-level operation. Such queries frequently combine nonlinear
field equations with Boolean structure, bit decompositions, and non-field
abstractions. Finite-field SMT provides a useful interface for these tasks:
verification tools construct logical properties while the solver manages
algebra and search. The cvc5 finite-field procedures and their public artifacts
established practical benchmarks for this setting~\cite{ff23,split24}.

We present \tool{}, an experimental Z3 branch~\cite{z3}, not an upstream
release. It integrates prime-field terms into the parser, APIs, models,
incremental contexts, and ground theory machinery. Native sorts remain available
inside uninterpreted functions, arrays, and datatypes, while structural
normalization and modular algebra avoid eager bit-vector expansion where
possible. Candidate field models guide shared-equality decisions; an exact
bit-vector bridge supplies fallback while preserving finite cardinality.
The implementation adds no CoCoA dependency.

Our contribution is an end-to-end tool integration, rather than a new complete
Gr\"obner-basis algorithm. We describe its semantic contract and combination
mechanism, distinguish explained conflicts from checkable proofs, and evaluate
all distinct finite-field inputs from three public artifacts. Category-level
coverage and failure plots reveal complementary strengths instead of treating
the pooled score as universal dominance. Real-project diagnostics and paired
optimization studies provide additional, separately identified evidence.
Independently checkable UNSAT certificates remain future work.

\section{Scope and User Interface}
\label{sec:scope}

For a prime $p$, the sort $\FF_p$ has domain $\{0,\ldots,p-1\}$ with arithmetic
modulo $p$. We support arbitrary-precision prime moduli, canonical numerals,
addition, multiplication, negation, equality, disequality, distinctness, Boolean
combinations, and field-valued conditionals. Each field has a separate sort;
terms of different moduli cannot be combined as though they belonged to one
field. The operator \flag{ff.bitsum} denotes $\sum_i 2^i x_i$ in the field and
does \emph{not} assert that any $x_i$ is Boolean. This distinction is essential
when reasoning about range constraints. The interface follows the prime-field
SMT-LIB conventions used by the evaluated artifacts; the broader finite-field
standardization effort is discussed by Hader and Ozdemir~\cite{ffa}.

The APIs expose field sorts, numerals, and arithmetic. A \qfff{} solver
selects the pure-field portfolio; general solvers also
support ground mixed queries under logic \flag{ALL}. Field terms retain their
sorts in foreign signatures, arrays, and datatype fields.

The following ground formula is UNSAT: a root of $x^2=1$ in $\FF_{17}$ is either
$1$ or $16$, and function congruence then contradicts one of the last assertions.
The integer-valued function is deliberately uninterpreted.

\begin{lstlisting}
(set-logic ALL)
(define-sort F () (_ FiniteField 17))
(declare-const x F)
(declare-fun tag (F) Int)
(assert (= (ff.mul x x) (as ff1 F)))
(assert (distinct (tag x) (tag (as ff1 F))))
(assert (distinct (tag x) (tag (as ff16 F))))
(check-sat)
\end{lstlisting}

Push/pop, assumptions, repeated checks, context translation, model evaluation,
and tracked cores are exercised by regression tests. These are operational
capabilities, not a claim of proof-producing support. General quantified field
solving and extension fields are outside scope. Universally stated properties
can nevertheless be checked by asking for a quantifier-free counterexample.
Moduli above $2^{64}$ undergo probable-prime screening, not certified primality
checking: the semantic contract requires the caller to supply a prime. Smaller
moduli use deterministic primality testing. Formal primality evidence and
algebraic proof reconstruction are deferred to a subsequent version.

\section{Solving Architecture}
\label{sec:architecture}

\paragraph{Native representation.}
A field declaration plugin connects parsing, rewriting, models, and APIs.
Sparse polynomials use canonical modular coefficients and ordered monomials.
Each modulus has a separate algebra problem. Foreign field-valued applications
are algebraic variables whose other semantics remain with their owning theory.

\paragraph{A staged portfolio.}
The default field strategy first applies \flag{ff-simplify}. It then attempts
\flag{ff-solve} for conjunctions of field literals, \flag{ff-sat} for lazy
Boolean/algebra search, and finally the native SMT theory plugin. Failure of a
bounded stage transfers the still-unresolved problem to the next stage; it is
not an UNSAT result. Field conditionals and eliminated wires have model
converters so returned assignments refer to the original vocabulary. The
standalone \flag{ff2bv} tactic is also available as an exact reference encoding.
Figure~\ref{fig:architecture} separates the pure-field strategies from the
mixed-theory path.

\begin{figure}[t]
\centering\small
\fbox{\begin{minipage}{.92\linewidth}
\centering Native AST and APIs $\to$ guarded preprocessing\\[4pt]
Conjunction algebra $\to$ lazy Boolean/algebra search\\[4pt]
\textit{Unresolved goals} $\downarrow$\\[4pt]
Field plugin $\leftrightarrow$ SAT/equalities $\leftrightarrow$ other theories\\[4pt]
\textit{Native algebra inconclusive} $\downarrow$\\[4pt]
Exact, range-constrained encode/decode bridge to BV
\end{minipage}}
\caption{Solving routes. Horizontal arrows between the first two algebra stages
mean portfolio fallback; the SMT arrows denote interaction during search.
Cancellation stops solving rather than restarting the portfolio.}
\label{fig:architecture}
\end{figure}

\paragraph{Bounded modular algebra.}
The engine performs normalization, variable elimination, critical-pair
processing, and univariate root reasoning. Product and completed-pair chain
criteria avoid redundant S-polynomials. Small primes below $2^{32}$ support
bounded F4-style batches using sparse modular row elimination; larger primes
use the arbitrary-precision sparse Buchberger path. These are bounded engines,
not a claim of the scalability of a mature computer-algebra backend. Default
local limits include two million algebra work units and 4,096 polynomial terms.
Sparse occurrence indices restrict substitution updates to affected constraints.

Disequalities can be represented algebraically by inverse witnesses:
$f\ne 0$ is equivalent over a field to $\exists t.\;ft=1$. Finding a nonzero
constant in the generated ideal establishes inconsistency. In contrast, merely
failing to derive a contradiction does not establish the existence of a point
in $\FF_p$. Univariate root extraction uses $\gcd(f,X^p-X)$ to restrict roots
to the intended prime field. Bounded sparse witness probes may discover SAT
assignments, but every accepted witness is checked against the original
polynomials and field AST. Failed probes cannot establish UNSAT.

\paragraph{Bounded root completion.}
A multivariate basis may describe only extension-field roots. When its leading
ideal contains a pure power of every active variable, a bounded pass reduces
$1,x,x^2,\ldots$ and searches for a linear dependence of their normal forms.
Applying the same row operations to formal powers yields a univariate ideal
consequence, which existing prime-field root reasoning can use. The pass stops
at degree 32 and charges at most the smaller of 30,000 work units and one
sixteenth of the remaining allowance. A failed search makes no logical claim.
The new default enables this pass together with diversified, fully checked SAT
witness probes; child probes do not recursively inherit diversification.

\paragraph{Boolean search and explanations.}
The lazy Boolean layer abstracts field atoms, asks the algebra engine about a
propositional assignment, and learns a clause from inconsistent premises.
Input dependencies are propagated through substitutions and reductions. These
sets support explained conflicts and cores; they omit polynomial multipliers
and therefore are not proof certificates. A bounded cache reuses exact basis
problems, checking the modulus, input polynomials, and premise dependencies.
This is exact reuse, not general incremental extension of a basis after adding
or removing arbitrary equations.

\section{Circuit Preprocessing and Representation Control}
\label{sec:preprocess}

Circuit expressions may have compact DAGs but large expanded polynomials.
The field rewriter folds constants, combines like terms, cancels additive terms
in equalities, and eliminates nonzero constant coefficients. Contextual
preprocessing propagates constants and acyclic wire definitions, retaining
original-model reconstruction and the premises of derived facts. Cyclic or
competing definitions remain constraints.

The equation $b(b-1)=0$ establishes Booleanity; it is never inferred merely
from a term's name or its occurrence in \flag{ff.bitsum}. After constant
propagation, suitable disjunctions expose Boolean domains through
$(u=0\lor v=0)\Leftrightarrow uv=0$. Bit-sum injectivity additionally requires
no modular wraparound. Signed interval propagation tests whether the range of
a Boolean linear combination contains an integer multiple of $p$, retaining
all equation, pin, and domain premises. Indicator recovery similarly requires
both defining equations of a zero test. Appendix~\ref{app:preprocess} gives
the conditions and the representation-control details.

Unconditional compaction can hide useful structure. The default instead retries
\flag{ff-solve} once with fresh defining variables only after polynomial-size
exhaustion during encoding. Replacing a subterm $t$ by fresh $v$ with $v=t$
preserves satisfiability and original-model reconstruction. Affine bit packs stay
visible. The retry receives at most one sixteenth of \flag{ff.max\_steps},
further capped by the first attempt's remaining budget; cancellation and
statistics cover both attempts. The rule tests no benchmark name, modulus, or
expected outcome. It also applies inside \flag{ff-sat}, but adds no independent
retry to the native SMT plugin. Other algebra and representation experiments
remain opt-in when measurements show no gain or regressions.

\section{Ground Theory Combination}
\label{sec:combination}

Finite fields are finite, so combination must not silently assume stable
infiniteness. The native theory plugin collects relevant equality classes and
assigned field equalities or disequalities, then solves a separate algebra
problem per modulus. It reconstructs values through the wire DAG and checks
all collected constraints. Equalities inherited from other theories are retained
as explicit premises, so an algebraic contradiction produces a conditional
lemma rather than an unexplained global assertion.

\paragraph{Shared model arrangements.}
Suppose two distinct shared roots receive the same candidate value in $\FF_p$.
The plugin asks the SAT/equality machinery to decide their equality. It does
\emph{not} assert that equality as an algebraic theorem: another valid field
model might separate them. The equality branch allows congruence, array
reasoning, and datatype reasoning to propagate; the disequality branch feeds a
new constraint into the next field check. Only roots observed by other theories
need these arrangements. Arranging every private circuit wire would introduce
irrelevant choices and destroy the benefit of DAG evaluation. Values are actual
field elements, so a candidate cannot create extra domain elements to avoid a
finite-cardinality conflict.

Conflicts preserve the exact SAT atoms used as premises. Rewriting a premise
into an equivalent but different atom can yield a learned clause already
satisfied by the current assignment, failing to reject the conflict. Dependencies
through wire normalization are consequently part of the implementation contract,
not optional bookkeeping. A bounded root pass also emits guarded product and
square clauses, such as $a^2=b^2\Rightarrow(a=b\lor a=-b)$, with characteristic
two handled without division by two.

\paragraph{Exact fallback with native signatures.}
When native algebra is inconclusive, the plugin enables a bridge for the current
SMT context. For $w=\lceil\log_2p\rceil$, private functions
$\mathit{enc}:\FF_p\to\mathrm{BV}_w$ and
$\mathit{dec}:\mathrm{BV}_w\to\FF_p$ enforce, on relevant terms,
\[
  \mathit{enc}(t)<p,\qquad
  \mathit{dec}(\mathit{enc}(t))=t.
\]
The bound is unnecessary when the bit-vector domain already equals the field
domain. Modular-operation constraints use widened arithmetic before reduction:
addition and multiplication must not lose information through machine-width
overflow. The left-inverse equation makes the representation injective; ordinary
function congruence transports equalities while foreign signatures keep their
field sorts. This construction provides an exact finite representation for the
ground field constraints, although the resulting BV problem can be expensive.
It does not make arbitrary mixed arithmetic theories decidable.

\paragraph{Incremental lifecycle.}
Backtracking invalidates candidate values and emitted-definition caches. Popped
axioms are re-emitted when surviving terms need them. Once selected, bridge mode
remains enabled for the context, so existing helper applications cannot later
be treated as unconstrained native variables. Helpers are hidden from exported
models. Cancellation returns an unresolved result instead of initiating fresh,
expensive work. Regression tests cover late field introduction, assumptions,
reset, context translation, and recovery after interruption.

\section{Evaluation}
\label{sec:evaluation}

We ask three questions: (R1) where does the integrated solver solve or miss
public benchmarks relative to cvc5; (R2) do the selected default changes improve
coverage without substantial aggregate cost; and (R3) what evidence supports
the semantic and mixed-theory integration? We distinguish complete corpus
coverage from development samples and application diagnostics.

\subsection{Corpus and Measurement Protocol}

The public artifacts of the CAV 2023 field solver, CAV 2024 split solver, and
FMCAD 2026 proof-production work provide 6,423 finite-field SMT files but only
\CorpusN{} byte-distinct inputs~\cite{artifact23,artifact24,artifact26}.
Identical files shared by artifacts are run once. Paper memberships overlap;
benchmark families in Figure~\ref{fig:categories} are disjoint. Alternate BV/NIA
encodings and CirC intermediate-format inputs are catalogued separately and
excluded from this finite-field SMT population. The artifact-specific
\flag{QF\_BVFF} logic name is replaced by \flag{ALL} for both solvers, without
changing assertions. For Z3, the unsupported option \flag{:incremental true}
 is removed; incremental commands and assertions remain unchanged.

Each primary run has a 10-second wall limit and a 4-GiB RSS limit sampled every
50 ms. Sampling permits memory overshoot. Fresh processes include startup,
parsing, solving, and shutdown. The machine is an Apple M2 Max with 12 cores
and 64 GiB RAM, running macOS; coverage measurements use eight workers without
CPU affinity. Z3 is a Release build without GMP. Both cvc5 configurations use
the same CoCoA-enabled 1.3.4 binary, with default GB or \flag{--ff-solver=split}.
Proof production is disabled for all measurements.

The historical Z3 columns use one frozen round-7 binary: 1,118 exact matches
are reused and 3,094 inputs are newly measured; seven Examples rows receive
one-worker compatibility reruns. cvc5 columns reuse the recorded 1.3.4 pass.
These are complete results for those binaries, not a contemporaneous latest-version
comparison. Historical timing and scheduling limit speedup claims. The paused
cvc5-main journal is excluded; 1.4.0 and main enter only correctness rechecks.
We do not reproduce the papers' original hardware or longer budgets.

The new-policy study uses 1,248 of those paper inputs plus eight separately
identified real-project diagnostics. Its two Z3 configurations are freshly
paired, with rotated run order and unchanged limits. Historical cvc5 rows are
joined by input identity for descriptive category comparisons; they are not
new paired measurements. The 120 hash-selected confirmation cases exclude the
previous cohort and development screen, but are not an untouched external corpus.

\subsection{Coverage, Complementarity, and Failure Categories}

\begin{figure}[t]
\centering\includegraphics[width=\linewidth]{round8-category-coverage.pdf}
\caption{Updated comparison on 1,248 selected paper inputs and eight separate
real-project queries (*). Every column uses the same inputs within a row.
cvc5 columns are historical 1.3.4; the Z3 columns are the new paired study.
Right: cases missed by the new policy but solved by cvc5 (red), or by none of
cvc5 GB, split, and the new policy (gray). Labels show red/total misses.
These subset counts do not replace the historical 4,212-input totals.}
\label{fig:categories}
\end{figure}

\begin{table}[t]
\caption{Historical full-corpus outcomes (4,212 inputs): cvc5 1.3.4 and
previous Z3 default. Other means unknown/error. Wrong
answers are excluded from solved counts. PAR-2 charges every unsuccessful run
20 seconds; lower is better.}
\label{tab:aggregate}
\centering\small\begin{tabular}{lrrrrrr}\toprule
Solver & Solved & Time & Mem. & Other & Wrong & PAR-2\\\midrule
\FailureRows
\bottomrule\end{tabular}
\end{table}

On the complete historical corpus, the previous default solves \ZthreeSolved{} inputs, compared with
\CvcSolved{} for GB and \SplitSolved{} for split. The total does not establish
universal dominance: \ReferenceOnly{} cases are solved by a cvc5 configuration
but missed by \tool{}, while \ZthreeOnly{} are solved only by \tool{} and
\NeitherSolved{} are unresolved by all three configurations. Figure~\ref{fig:categories}
shows the newly measured subset separately, including both coverage gains.
The historical full table, category view, and per-input status map are in
Appendix~\ref{app:data}.

The historical gaps concentrate in Small (34), TV (9), and TV-pureFF (3).
The new paired study improves Examples from 9/10 to 10/10 and QED2 from 71/100
to 72/100. On its 1,248 paper inputs, the new policy solves 885, versus 466 GB
and 713 split successes in the historical references; 29 cvc5-union-only cases
remain. This selected subset cannot establish a new whole-corpus ranking.

\begin{figure}[t]
\centering\includegraphics[width=\linewidth]{round8-paired-runtime.pdf}
\caption{Fresh paired Z3 measurements on all 1,256 study inputs. Left: successful
runs only. Right: the 886 common successes and two new solves; previous timeouts
are placed at the T/O boundary, not assigned a successful runtime. Points below
the diagonal favor the new policy. Both configurations fail on 368 other inputs.}
\label{fig:timing}
\end{figure}

\subsection{Controlled Optimization Studies}

The new default enables \flag{ff.model\_search}, combining the bounded
univariate-consequence pass with diversified SAT witness probes. On 1,256 inputs,
coverage rises from 886 to 888, with two gains, no losses, and no contradictory
answers. The previous cohort plus all Examples gives 800 to 802 successes on
1,136 inputs; the separate 120-input confirmation cohort remains at 86. The
confirmation cohort supplies regression evidence, not a new coverage gain.
Summed time on 886 common successes changes from 227.992 to 229.028 seconds
(+0.45\%); the geometric candidate/baseline ratio for common cases taking at
least 50 ms in either run is 0.996. Figure~\ref{fig:timing} shows the paired data.

Three isolated repetitions confirm both gains: the cyclic Example changes from
timeout to UNSAT (median 0.018 s), and Montgomery2Edwards from timeout to SAT
(0.026 s). MACI Merkle inclusion remains at 4.681 versus 4.684 s. Two short SAT
cases regress from 0.036 to 0.125 s and 0.042 to 0.168 s. Thus the coverage gain
has comparable aggregate cost, but is not uniform acceleration.

The broad run uses a frozen experimental binary with the option explicitly on;
the retained build changes its public default to true. A 27-input replay checks
agreement between these configurations, including all Examples and 12 Poseidon
sentinels. Root-only completion is independently available. Quotient Frobenius
facts and adaptive matrix admission remain opt-in: they add no solved cases in
their measured screens. Appendix~\ref{app:round8} reports ablations, resource
failures, validation, and exact binary identities.

Earlier representation-control and bit-domain studies are reported separately
in Appendix~\ref{app:preprocess}; their overlapping samples are not pooled.

\subsection{Real Circuits and Mixed-Theory Workloads}

Eight separate queries use published Picus circuits from Dark Forest, Hermez,
MACI, and a Succinct Plonky2 verifier~\cite{picus}. Two complete circuit copies
share inputs and differ on outputs; UNSAT establishes output determinism of the
encoding, not functional correctness or security. This monolithic test does not
reproduce Picus's decomposition pipeline. All use BN254, including emulated
Plonky2 arithmetic. Both paired Z3 policies solve 3/8; historical cvc5 solves
0/8. The five unresolved cases limit application-wide conclusions.

A separate historical integration experiment covers 16 generated UF/array
and Poseidon queries over BN254 and BLS12-381. Native SMT and the default Z3
path solve all 16; each cvc5 backend solves eight. Different binaries and
3-second internal/5-second external limits prevent pooling these results with
the main corpus. The latest 12 Poseidon sentinels preserve Z3 coverage and
validated SAT models, without a new three-solver timing claim.

\section{Correctness Evidence and Limitations}
\label{sec:limits}

The retained build passes 16 regression commands covering parser/API behavior,
incrementality, exhaustive small-field semantics, model evaluation, cores,
resource limits, and proof-request rejection. Mixed-theory tests enumerate 105
queries and unary-function tables through default, native, persistent, and
forced-fallback routes. Algebra tests perform 512 basis, S-pair, and dependency
checks over 64 random systems, including primes near $2^{32}$ and BN254.
The new root and matrix tests are detailed in Appendix~\ref{app:round8}.

All 1,295 historical round-7 single-query SAT results and all 559 new paired-study
SAT model replays pass independent original-assertion evaluation. These are
overlapping validation populations, not extra benchmark inputs. Representative
final-build replays also validate (Appendix~\ref{app:round8}). No independent
UNSAT certificate is produced.

One recorded cvc5 split answer in the public corpus was refuted by independently
validated SAT assignments. The raw answer remains archived but is scored as
wrong, not successful. Targeted rechecks reproduce it on cvc5 1.4.0 and main at commit
\texttt{72f647eb75c0}, three fresh runs per backend and version. This
single finding is neither a general correctness comparison nor a justification
for treating all Z3 UNSAT answers as proved. All other definite cross-solver
answers are checked for agreement in the joined results.

The principal threats to validity are related benchmark generators, adaptive
development, a short time limit, and one measurement platform. Hash-selected
confirmation sets reduce direct tuning but are not independent generators.
The full pass is a descriptive evaluation, not an untouched randomized trial.
Historical cvc5 timing, resource failures, unresolved basis growth, and the
absence of a full-corpus run of the new default limit performance conclusions. A current-version comparison, Linux validation, and
an end-to-end Blaster application remain important submission work.

Proof production is unsupported. A future certificate layer must record
normalization identities, reduction multipliers, elimination substitutions,
inverse-witness introductions, root coverage, quotient normal-form identities, guarded bit deductions, and
the Boolean derivation. Fresh compact definitions and the combination bridge also
need checked extension arguments. Model validation and dependency sets do not
replace these obligations. The integration is not yet a proof-reconstructing
Lean backend, and general quantified solving remains outside its acceptance
claims.

\section{Related Work and Conclusion}

Ozdemir et al. introduced practical prime-field reasoning in cvc5 using
algebraic methods and explained conflicts~\cite{ff23}. Their split-basis solver
uses cooperating bases and specialized propagation~\cite{split24}; our bounded
split prepass does not reimplement it. Hader et al.'s MCSat finite-field plugin
in Yices2 uses zero decomposition and finite conflict explanations~\cite{yices}.
We have not benchmarked Yices2 and make no overall state-of-the-art claim.

Recent work adds proof production and checking in Pacheck and Lean for cvc5
finite fields~\cite{proof26}, a capability absent here. Our contribution is
native integration into Z3's APIs and ground theory infrastructure, with
circuit-preserving normalization, validated models, and exact fallback.
The category-level evaluation identifies useful coverage and remaining gaps.

\paragraph{Data availability.}
The local artifact includes source/binary identities, inputs, results,
validation, figures, and editable LaTeX. A permanent public release is pending.

\clearpage
\bibliographystyle{splncs04}
\bibliography{references}

\clearpage
\appendix
\section{Detailed Results and Benchmark Inventory}\label{app:data}

\begin{table}[h]
\caption{Historical complete paper-corpus coverage (round-7 Z3, cvc5 1.3.4). ``cvc-only'' counts the cvc5 union's
successes missed by Z3+FF, not a separate solver configuration.}
\centering\small\begin{tabular}{lrrrrr}\toprule
Family & Inputs & GB & Split & Z3+FF & cvc-only\\\midrule
\FamilyRows
\bottomrule\end{tabular}
\end{table}

TV is compiler translation validation and TV-pureFF its pure-field encoding.
CirC-D is operator determinism; CirC-S is mixed BV/field operator soundness.
QED2 contains circomlib determinism queries. Seq contains bit-sum sequences;
Small contains small-field polynomial systems; ASHR contains arithmetic-shift
queries. Examples are illustrative artifact inputs not already assigned to an
experimental family. Duplicate examples inherit the experimental category.

\begin{table}[h]
\caption{Historical artifact and modulus views (round-7 Z3, cvc5 1.3.4). Artifact rows overlap; field-size rows
partition inputs by the largest field modulus.}
\centering\small\begin{tabular}{lrrrr}\toprule
Selection & Inputs & GB & Split & Z3+FF\\\midrule
\PaperRows
\midrule
\SizeRows
\bottomrule\end{tabular}
\end{table}

\begin{figure}[p]
\centering\includegraphics[width=\linewidth]{all-input-statuses-paper.pdf}
\caption{Historical full-corpus status map (round-7 Z3, cvc5 1.3.4).
Every primary input is a column within its category, with three solver
rows. Columns are ordered by outcome pattern and then input hash; they do not
represent chronological execution. Dense columns are easier to inspect in the
standalone figure or searchable HTML/CSV exports. SAT and UNSAT are both solved;
wrong, timeout, memory, unknown, and error are distinct failure outcomes.}
\end{figure}

\begin{figure}[p]
\centering\includegraphics[width=\linewidth]{category-coverage.pdf}
\caption{Historical full-corpus category comparison: all 4,212 inputs, cvc5
1.3.4 and the previous (round-7) Z3 default. The new default has no full-corpus
measurement yet; Figure~\ref{fig:categories} is its measured subset.}
\end{figure}
\begin{figure}[p]
\centering\includegraphics[width=\linewidth]{runtime-overview.pdf}
\caption{Historical full-corpus cactus and performance profile. Different
measurement dates and scheduling preclude a controlled speedup interpretation.}
\end{figure}

\clearpage
\section{Supplementary Evidence and Reproduction}

\begin{table}[h]
\caption{Separate real-project output-determinism queries, at 10 seconds.
Numbers are successful UNSAT runtimes in seconds; T is timeout and M is the
sampled memory limit. cvc5 timings use the earlier one-worker experiment;
Z3 rows in this historical table use the round-7 eight-worker pass.}
\centering\scriptsize\begin{tabular}{lrrr}\toprule
Query & GB & Split & Z3+FF\\\midrule
\RealRows
\bottomrule\end{tabular}
\end{table}

The evaluation populations serve different purposes and are not added together:
(i) 4,212 distinct paper inputs provide full public-artifact coverage;
(ii) eight real-project queries are application diagnostics;
(iii) 16 generated mixed-theory workloads are historical integration evidence;
(iv) 12 current Poseidon sentinels are regression checks; and
(v) the overlapping optimization cohorts support paired default-selection
experiments. The remaining API, random-system, root, rewrite, and combination
suites establish semantic checks rather than a competitive benchmark score.
The historical full-corpus Z3 result is not spliced into the new-policy subset.

\paragraph{Rebuilding the evaluation.}
The corpus manifest records original artifact membership and SHA-256 identities.
The frozen Z3 executable is identified by the prefix \texttt{14772994298957b7};
the cvc5 reference by \texttt{802ca3d002b4b61e}. Full hashes and commands are in
the machine-readable metadata. The Z3 source overlay applies to the recorded
base commit; it includes the field implementation and tests. Fresh processes
must receive the same logic-name adaptation and limits. The current completion
journal identifies every reused row and every newly measured row. It must not
be interpreted as three contemporaneously interleaved solver runs.

The figure generator derives tables and plots from those journals, checks that
all primary inputs have one round-7 Z3 row, applies the recorded cvc5 wrong-answer
adjudication, and rejects unresolved definite-answer disagreements. Raw solver
results, unsuccessful experiments, and independent SAT-validation records remain
available. Rebuilding the document requires the bundled LNCS class, standard
LaTeX packages, BibTeX, and the generated figures. Author and institution metadata
are supplied by the authors; this working draft has not been submitted.

\clearpage
\section{Preprocessing Conditions and Experimental Algorithms}
\label{app:preprocess}

Preprocessing matters because a circuit often has a compact expression DAG
but a very large expanded polynomial. The global rewriter folds constants,
removes zero and unit factors, combines like terms, cancels additive terms in
equalities, and eliminates nonzero constant coefficients. The contextual pass
propagates constants and substitutes acyclic wire definitions while retaining
a model-reconstruction map. Competing or cyclic definitions remain constraints.
The mixed-theory path similarly evaluates wire definitions through a shared DAG.

\paragraph{Guarded indicator recovery.}
A zero test is recognized only from its defining constraints. For example,
$t z=0$ and $t w+z=1$ imply $z=1$ when $t=0$, and $z=0$ when $t\ne0$.
Both premises are required: keeping only the annihilation equation would permit
spurious values when $t=0$. Recovered facts retain their dependencies, including
substituted definitions. This preserves the distinction between an algebraic
consequence and an exporter convention.

\paragraph{Boolean and bit-sum reasoning.}
The equation $b(b-1)=0$ implies $b\in\{0,1\}$ in a field. Disjunctive Boolean
domains can be exposed after constant propagation using
$(u=0\lor v=0)\Leftrightarrow uv=0$. Matching bit sums additionally requires
Booleanity of every digit. For no-wrap injectivity, $2^n\le p$ ensures that an
$n$-bit integer has a unique residue. Signed interval propagation generalizes
this check: bound a linear combination of Boolean digits and test whether its
integer interval contains a multiple of $p$. Testing both values of a digit can
establish a pin or a contradiction. The implementation retains the equation,
used pins, and Boolean-domain premises. It neither assumes Booleanity from the
name of a term nor ignores modular wraparound.

\paragraph{A compact-encoding retry.}
Unconditional fresh-variable introduction can obscure useful structure and
regress performance. Instead, the default preserves the original representation
first. If polynomial encoding, specifically, exhausts its size bound,
\flag{ff-solve} retries once using exact fresh definitions. For a subterm $t$,
replacing an occurrence by fresh $v$ and adding $v=t$ preserves satisfiability
and supports reconstruction of the original model. Affine bit packs remain
visible, and nonlinear definitions are protected against immediate elimination
that would recreate the expansion. The retry receives at most one sixteenth of
\flag{ff.max\_steps}, further capped by the work remaining after the first
encoding. Cancellation and statistics cover both attempts. This tactic-level
retry is also used by \flag{ff-sat}, but is not independently inserted into the
native SMT plugin.

The budget is a general bound, without tests on benchmark identity, modulus, or
expected status. Other implemented experiments---sugar scheduling,
Gebauer--M\"oller pair filtering, divisibility support masks, geobuckets,
coefficient shortcuts, scalar recovery from matrix storage exhaustion, and
substitution-growth estimates---remain disabled by default. Their correctness
checks passed, but the measured screens did not justify enabling them. We
report this negative evidence instead of treating every implementation as a
performance contribution.

\paragraph{Default-selection evidence.}
This retained result followed a rejected version: giving compact encoding the
entire remaining allowance caused a reproducible MACI slowdown from a 4.615-s
median to 6.647 s. With the retry capped to one sixteenth of the local work
budget, final isolated medians are 4.699 and 4.846 s. Two short determinism
queries still incur roughly 0.1 s overhead. Both gained cases repeat across
three isolated runs, with candidate medians of 3.082 s (UNSAT) and 0.405 s (SAT),
against baseline timeouts. The tradeoff is a small coverage improvement at
comparable aggregate time, not a claim of uniform acceleration.

The preceding bit-domain study improved from 563 to 581 solves on a separate
827-input paired sample, with 18 gains and no losses. Twelve gains occurred in
its 179-input fresh confirmation cohort. These samples overlap across rounds
and must not be pooled as independent trials. Direct compaction, split prepasses,
and several algebraic heuristics were left disabled when their screens showed
no benefit or regressions. This measurement policy avoids enabling a technique
solely because it is plausible in another solver architecture.


\clearpage
\section{New-Policy Study: Ablations and Remaining Gaps}\label{app:round8}

\begin{table}[h]
\caption{Updated subset coverage at 10 seconds. Old/new are the paired Z3
policies; GB and split are historical cvc5 1.3.4 results on the same inputs.
The final column counts cvc5-union successes missed by the new policy.}
\centering\small\begin{tabular}{lrrrrrr}\toprule
Family & Inputs & GB & Split & Old & New & cvc-only\\\midrule
\RoundEightRows
\bottomrule\end{tabular}
\end{table}

\begin{figure}[h]
\centering\includegraphics[width=\linewidth]{round8-reference-runtime.pdf}
\caption{Descriptive timing on the 1,248 selected paper inputs, excluding the
eight real-project diagnostics. cvc5 timings are historical; these plots are
not a contemporaneous three-solver experiment. Failed runs are absent from the
cactus and have infinite performance-profile ratios.}
\end{figure}

\paragraph{Candidate separation.}
The initial 113-input screen solves 51 with the previous default, 52 with
root-only completion, 51 with quotient-field facts, 52 with both, and 53 with
model search. Root completion supplies the cyclic UNSAT gain; diversified
witness probes additionally supply Montgomery2Edwards SAT. No definite answers
conflict. This exploratory screen overlapped development work and cannot by
itself support performance equivalence; the clean paired study does.

Quotient-field reasoning checks that the leading ideal bounds every active
variable and enumerates at most 128 standard monomials. It computes reduced
Frobenius facts $\mathrm{NF}(x^p-x)$ using repeated squaring and reduction, then
rebuilds a bounded basis. Such facts restrict solutions to the intended field;
resource exhaustion is not inconsistency. The candidate remains off because
its screen supplies no additional solved inputs.

Adaptive matrix admission replaces the fixed 1,024 symbolic-reducer cap by
conservative accounting against 16 MiB of symbolic storage, retaining the other
work, column, term, and cancellation bounds. This is not a total-process memory
bound. Alone it solves the same 51/113 cases and increases common-solve time
by 5.3\%; combined with model search it remains at 53. It therefore stays off.

\paragraph{Where additional root reasoning does not suffice.}
All 34 Small-family gaps in the historical cvc5-union-only set exhaust an
algebra budget before root completion: 18 work, 15 matrix, and one term limit.
These are three-second direct-algebra diagnostics, not measurements of the
full ten-second portfolio. A ten-second follow-up on those 34 cases compares
the old matrix policy at eight million work units with adaptive admission at
eight and 32 million. Each solves the same one case. Additional diagnostics
confirm that the adaptive path admits extra reducers but still exhausts matrix
resources. These measurements point to basis construction and sparse matrix
scalability as the larger remaining bottleneck in this family.

\begin{figure}[t]
\centering\includegraphics[width=\linewidth]{round8-isolated.pdf}
\caption{Three isolated repetitions per configuration. Gains survive isolation;
MACI is unchanged, while two short SAT cases become slower. Timeout bars mark
the ten-second limit and are not successful runtimes.}
\end{figure}

\paragraph{Semantic checks.}
Tests include an exact F5 univariate ideal consequence, independent scalar
ideal membership, exhaustive premise-necessity checks, an F3 system with only
extension-field solutions, cancellation, and matrix tests above 1,024 reducers.
Random quadratic systems and root families validate models and explanations.
All 16 regression commands pass. All 559 primary SAT model replays satisfy the
original assertions; 27 final-build comparisons preserve complete answer
sequences, including multi-query Examples. Independent model checks cannot
certify UNSAT; algebraic proof reconstruction remains future work.

\paragraph{Excluded measurements and runner reliability.}
An earlier wide run was interrupted after host disturbance and a supervision
failure that left solver children running. It is archived but excluded from
acceptance. The repaired runner owns worker/solver process groups, reaps them
on outer failures, journals infrastructure errors separately, and rejects
unsafe mixed-log resumes. Five focused runner tests pass. The clean 2,512-run
study has zero infrastructure errors. Isolated reruns remove both apparent
coverage losses from the interrupted run.

\paragraph{Reproduction and binary identity.}
The previous binary is identified by SHA-256 prefix \texttt{14772994298957b7}.
The broad candidate measurements use experimental binary
\texttt{02b948795788f98f} with \flag{smt.ff.model\_search=true}; the retained
binary \texttt{4c63ed6003400663} makes that policy the public default. Full hashes,
source overlays, exact inputs, paired logs, isolated runs, model validation,
and regression results accompany the artifact. Other new experimental options
remain false. Child-engine model search remains false to bound nested probes.

The updated figure generator joins exactly the measured 1,248 paper identities
and eight real-query identities, verifies definite-answer agreement, and keeps
historical cvc5 timing distinct. It never fills unmeasured current-default
inputs with old results. The current cvc5-main campaign was paused after
1,018 of 12,636 planned runs and is excluded from this paper's aggregate tables.
A completed contemporary full-corpus comparison remains necessary before a
stronger current-version performance claim.

\end{document}
